\documentclass[12pt]{article}
\usepackage[a4paper,margin=1in]{geometry}\usepackage[T1]{fontenc}
\usepackage{lmodern,microtype}\usepackage{amsmath,amssymb,amsthm,mathtools}
\usepackage{booktabs,array,enumitem,xcolor}\usepackage[numbers,sort&compress]{natbib}
\usepackage[colorlinks=true,linkcolor=blue!55!black,citecolor=blue!55!black,urlcolor=blue!55!black]{hyperref}
\linespread{1.07}
\newtheorem{theorem}{Theorem}[section]\newtheorem{proposition}[theorem]{Proposition}
\newtheorem{lemma}[theorem]{Lemma}\newtheorem{corollary}[theorem]{Corollary}
\theoremstyle{remark}\newtheorem{remark}[theorem]{Remark}
\newcommand{\norm}[1]{\left\lVert#1\right\rVert}

\title{Boundary Stratification of the Centered Quartet Manifold\\
\large Discriminant Factorization and Uniform Collapse to a Double Pair}
\author{Liang Wang}\date{July 2026}
\begin{document}\maketitle
\begin{abstract}
The finite shape clock of RH-214 moves monotonically in the axial coordinate
$u$, and RH-215 finds a power-gap model to be the best of three finite
extrapolants.  Neither result proves a limit.  This paper instead answers the
conditional geometric question exactly: what happens to a centered
conjugate quartet if $u$ approaches the boundary value one?

For the canonical shape polynomial $Q_{u,\eta}$, we prove
\[
 \operatorname{Disc}Q_{u,\eta}
 =256(1-u)^2(1-\eta^2)
   \bigl[4u+(1-u)^2\eta^2\bigr]^2.
\]
Hence degeneracy occurs precisely on the axial edge $u=1$, on either edge
$|\eta|=1$, or at the isolated point $(u,\eta)=(0,0)$.  The strata correspond
respectively to two real double roots, one real double pair, and coincident
imaginary conjugate pairs.

Uniformly for $|\eta|\le1$, the root multiset approaches
$\{1,1,-1,-1\}$ as $u\to1$, with an explicit matching bound
\[
 d_{\rm root}\le \sqrt{2(1-u)}+
 \frac{1-u}{1+\sqrt u}.
\]
The coefficient distance to $(z^2-1)^2$ is exactly $4(1-u)$ in maximum norm.
Thus axial convergence, coefficient convergence, and root-multiset
convergence are equivalent.  Six hundred random checks verify the formulas
to roundoff.  The physical data move toward this stratum but reach only
$u\approx0.718$ at the finest anchor; no physical convergence claim is made.
\end{abstract}

\section{Conditional geometry without a convergence claim}

RH-213 gives the exact canonical roots
\begin{equation}\label{eq:roots}
 a\pm ib,\qquad -a\pm id,
\end{equation}
where
\begin{equation}\label{eq:abd}
 a=\sqrt u,qquad
 b^2=(1-u)(1+\eta),qquad
 d^2=(1-u)(1-\eta)
\end{equation}
for $(u,\eta)\in[0,1]\times[-1,1]$
\citep{WangRH213}.  The polynomial is
\begin{equation}\label{eq:poly}
 Q_{u,\eta}(z)=z^4+(2-4u)z^2
 +4\sqrt u(1-u)\eta z+1-\eta^2(1-u)^2.
\end{equation}

RH-214 observes finite monotonicity of $u$ and RH-215 tests three candidate
tails \citep{WangRH214,WangRH215}.  We do not assume any of those tails.  All
results below hold for an arbitrary point of the compact shape rectangle.

\section{Quadratic factorization}

Write
\begin{equation}\label{eq:factors}
 Q_+(z)=(z-a)^2+b^2,qquad
 Q_-(z)=(z+a)^2+d^2,qquad Q=Q_+Q_-.
\end{equation}
The discriminants of the two monic quadratics are $-4b^2$ and $-4d^2$.

\begin{lemma}[Cross resultant]\label{lem:resultant}
The resultant of the two factors is
\begin{equation}\label{eq:resultant}
 \operatorname{Res}(Q_+,Q_-)
 =4\bigl[4u+(1-u)^2\eta^2\bigr].
\end{equation}
\end{lemma}
\begin{proof}
Evaluating $Q_-$ at the two roots $a\pm ib$ of $Q_+$ gives
\begin{align}
 \operatorname{Res}(Q_+,Q_-)
 &=\bigl((2a+ib)^2+d^2\bigr)
   \bigl((2a-ib)^2+d^2\bigr)\\
 &=\bigl(4a^2-b^2+d^2\bigr)^2+16a^2b^2.
\end{align}
Using \eqref{eq:abd}, the terms linear in $\eta$ cancel and the expression
reduces to \eqref{eq:resultant}.
\end{proof}

\section{Complete discriminant factorization}

For monic polynomials,
\begin{equation}
 \operatorname{Disc}(Q_+Q_-)
 =\operatorname{Disc}(Q_+)\operatorname{Disc}(Q_-)
  \operatorname{Res}(Q_+,Q_-)^2.
\end{equation}

\begin{theorem}[Shape discriminant]\label{thm:disc}
For the canonical quartet,
\begin{equation}\label{eq:disc}
 \boxed{
 \operatorname{Disc}Q_{u,\eta}
 =256(1-u)^2(1-\eta^2)
 \bigl[4u+(1-u)^2\eta^2\bigr]^2.}
\end{equation}
\end{theorem}
\begin{proof}
Since $b^2d^2=(1-u)^2(1-\eta^2)$, the product of the two quadratic
discriminants is $16(1-u)^2(1-\eta^2)$.  Multiplying by the square of
\eqref{eq:resultant} proves \eqref{eq:disc}.
\end{proof}

\section{Degeneracy strata}

All factors in \eqref{eq:disc} are nonnegative on the shape rectangle.

\begin{corollary}[Complete zero set]\label{cor:strata}
The quartet has a repeated root if and only if at least one of the following
holds:
\begin{enumerate}[leftmargin=2em]
\item $u=1$;
\item $|\eta|=1$;
\item $(u,\eta)=(0,0)$.
\end{enumerate}
\end{corollary}
\begin{proof}
The first two factors in \eqref{eq:disc} give the first two cases.  The final
square bracket vanishes only when both nonnegative terms $4u$ and
$(1-u)^2\eta^2$ vanish, which is exactly $u=0$, $\eta=0$.
\end{proof}

The geometry of each stratum is explicit:
\begin{description}[leftmargin=3.2em,style=nextline]
\item[$u=1$]
$a=1$ and $b=d=0$, giving roots $1,1,-1,-1$.
\item[$\eta=1$]
$d=0$, so the negative-real pair coalesces at $-\sqrt u$; for $\eta=-1$,
the positive-real pair coalesces at $+\sqrt u$.
\item[$(u,\eta)=(0,0)$]
$a=0$ and $b=d=1$, so both quadratic factors equal $z^2+1$ and the roots
$\pm i$ each have multiplicity two.
\end{description}

Intersections carry the corresponding higher overlap, but introduce no new
zero set.

\section{Uniform collapse along the axial boundary}

Let $\mathcal B=\{1,1,-1,-1\}$, with multiplicity, and let
$d_{\rm match}$ be the minimum over root permutations of the maximum paired
distance.

\begin{theorem}[Uniform root-collapse estimate]\label{thm:collapse}
For every $(u,\eta)$ in the shape rectangle,
\begin{equation}\label{eq:rootbound}
 d_{\rm match}(\operatorname{roots}Q_{u,\eta},\mathcal B)
 \le \sqrt{2(1-u)}+\frac{1-u}{1+\sqrt u}.
\end{equation}
In particular the convergence to $\mathcal B$ is uniform in $\eta$ as
$u\to1$.
\end{theorem}
\begin{proof}
Match the two roots with real part $a=\sqrt u$ to $1$ and the two with real
part $-a$ to $-1$.  For the positive pair,
\begin{equation}
 |a\pm ib-1|\le |1-a|+b
 \le \frac{1-u}{1+\sqrt u}+\sqrt{2(1-u)},
\end{equation}
because $b^2\le2(1-u)$.  The same estimate holds for the negative pair using
$d^2\le2(1-u)$.
\end{proof}

The square-root term is natural at root level: imaginary heights vanish like
$\sqrt{1-u}$ in the worst transverse direction.

\section{Coefficient collapse and equivalence}

The boundary polynomial is
\begin{equation}
 Q_{1,\eta}(z)=(z^2-1)^2=z^4-2z^2+1,
\end{equation}
independent of $\eta$.

\begin{proposition}[Exact coefficient distance]\label{prop:coeffdistance}
Let $C(u,\eta)=(1,0,c_2,c_3,c_4)$ and
$C_\star=(1,0,-2,0,1)$.  Then
\begin{equation}\label{eq:coeffdistance}
 \norm{C(u,\eta)-C_\star}_\infty=4(1-u).
\end{equation}
\end{proposition}
\begin{proof}
The $c_2$ difference is exactly $4(1-u)$.  Moreover
\begin{equation}
 |c_3|\le4\sqrt u(1-u)\le4(1-u),qquad
 |c_4-1|\le(1-u)^2\le4(1-u).
\end{equation}
Thus no other coefficient exceeds the $c_2$ difference.
\end{proof}

\begin{corollary}[Equivalent axial convergence]\label{cor:equiv}
For any sequence $(u_n,\eta_n)$ in the shape rectangle, the following are
equivalent:
\begin{enumerate}[leftmargin=2em]
\item $u_n\to1$;
\item $Q_{u_n,\eta_n}\to(z^2-1)^2$ coefficientwise;
\item the root multisets converge to $\mathcal B$ in matching distance.
\end{enumerate}
\end{corollary}
\begin{proof}
Proposition~\ref{prop:coeffdistance} proves equivalence of the first two.
Theorem~\ref{thm:collapse} gives the first implies the third.  Conversely,
root convergence forces the positive real branch coordinate
$\sqrt{u_n}$ to approach one.
\end{proof}

\section{Machine audit and physical position}

Six hundred uniformly sampled parameter pairs verify
Theorem~\ref{thm:disc} by direct root-product discriminants.  The maximum
relative error, normalized by $\max(1,|\operatorname{Disc}|)$, is below
$10^{-14}$; no violation of \eqref{eq:rootbound} occurs.

For the physical atlas, $u$ rises from about $0.02$ to about $0.718$.
Consequently the maximum-norm coefficient distance in
\eqref{eq:coeffdistance} falls monotonically, but at the finest level it is
still approximately $1.13$.  The matching distance to the double pair is
about $0.577$.  These are directional observations, not evidence that the
asymptotic regime has been reached.

\section{Route consequence and claim boundary}

The axial boundary is now mathematically understood.  If a later analytic or
validated argument proves $u_\sigma\to1$, no separate control of $\eta$ is
needed for coefficient convergence: transverse dependence is uniformly
suppressed.  RH-217 quantifies that suppression differentially.

Even a proved fixed-quartet limit would still have only four roots, with two
limiting locations.  It cannot by itself provide a growing spectral divisor.
No physical limit, determinant construction, Gate-A closure, Hilbert--P\'olya
operator, or zeta statement is asserted here.

\bibliographystyle{plainnat}\bibliography{references}\end{document}
